% Technion Physics Dept training laboratories report template
%
% Based on "Report Preparation Guidelines" by Gil Ben-Ari and Dr.Yulia Preezant
% 114081 - Physics Laboratory 1 - International School
%
% Adapted to LaTeX by Evgeny Kolonsky 2023

% Preamble
\documentclass[a4paper,12pt]{article}

% Packages for figures, tables and their captions
\usepackage{graphicx}
\usepackage{caption}

% Package for SI units
% Documentation: https://texdoc.org/serve/siunitx/0
\usepackage{siunitx}
\sisetup{uncertainty-mode=separate} % Like 1.5 +/- 0.1

% Package for neat equations 
%\usepackage{amsmath} 

% Package for neat code listings
% Documentation: https://www.overleaf.com/learn/latex/Code_listing
\usepackage{listings}
\input{listings_settings.tex}


% Packages for citation and referencing
\usepackage[sorting=none]{biblatex} % sort by citation order
\addbibresource{bibliography.bib}  % Replace with the actual path to your .bib file

\begin{document}

% Title page

\begin{titlepage}
  \centering
  Technion Physics Dept, Laboratory 1 - Mechanics. \\
  Laboratory Experiment Report

  \vspace{2cm}
  \begin{tabular}{@{}l@{\hspace{1cm}}l@{}}
  Course No.:& 114081 \\
  Term: & Summer semester 2023\\
  Group  No.:& 12 \\
  Instructor's name: & Gabi Weingart \\
  Workplace No.: & 6 \\
  Experiment: & Simple Pendulum\\
  \\
  Prepared by: & Moshe Cohen, \\
               & ID. 123456789 \\
               & \texttt{moshec@campus.technion.ac.il} \\
               \\
               & Orian Levi, \\
               & ID 3456789123\\
               & \texttt{orianl@campus.technion.ac.il} \\
  \\
  Date: & {01 August 2023}\\
  \end{tabular}


  \vspace*{\fill}
  
  
\end{titlepage}

\pagenumbering{arabic} 

\begin{center}
{\LARGE \textbf{Short and Substantial Title}} \\  % Title
\small{The report title describes the essence of the experiment. The title should be significant (not formal) and particular (not general).} 
\end{center}


\begin{abstract}
\emph{Important!} This section does not coincide with the "Lab Objective" section, which appears at the beginning of each chapter in the laboratory manual. This section is brief (one or two paragraphs) and is written as a narrative as opposed to a list. 

This section answers the following questions: What is the physical phenomenon? What is being measured? What physical quantities are derived from the measurements? 
\end{abstract}

\section{Introduction}

\emph{Important!} This section contains only information relevant to the goals of the experiment. 

This section is short, up to one-half page.

This section contains a brief description of the theoretical model used to predict the experimental results using 1 \textendash \  5 equations, all of which are numbered. Include a description of the model assumptions and limitations relevant to the experiment.

For each equation, specify its parameters and physical units. Example:
\begin{equation} \label{eq1}
\vec{F} = m \vec{a},
\end{equation}
where $\vec{F}$ is force vector in \unit{\newton}, $m$ is mass in \unit{\kilogram} and $\vec{a}$ is acceleration vector in \unit{\meter\per\s\squared}.

This section contains an illustration of the experimental system and the experimental conditions. As necessary, include a description of the measuring instruments and their measurement accuracies. 

To understand better the basics of computer typesetting refer to \TeX \ book \cite{KhuthTeX}.

\section{Experimental results}

This section is written as a sequential and logical story describing the execution of the experiment, from initial calibration, through the processing of the data to the calculation of the desired outcome and accompanying estimate of experimental errors.  The section may contain graphs, schemes, tables, and any visual object to represent the results.  The text, graphs, and tables are arranged in a logical order, which begins from the results of the raw measurement to the discussion of the conclusions and explanations that connect them. This section may be subdivided with descriptive subheadings as an aid in developing the narrative as necessary.

Example of uncertainty representation: 
$$
g = \qty{9.7949(1)}{\meter\per\sec\squared}.
$$

The experiment results should be represented with uncertainties. To understand better uncertainties and errors in experiments refer to \cite{Tailor01}.


\subsection{Requirements for graphs}
The graph is presented clearly. The font size is 10-12 point. 
This axis is appropriately sized to display the data clearly on the graph. No gray backgrounds.
Measured data is presented as discrete points (without connecting lines). The results of theoretical calculations are displayed as a continuous curve.
Display the numerical labels of the axes outside the graph alongside each axis.
Use legends when there is more than one data type displayed on the graph.
Graph captions are required.  Place them immediately below the graph, and use them to describe what the graph represents. There is no need to add a title to the graph because there is an explanation in the caption below.
Write the axis titles and units. The axes are explained in the caption of the graph, if necessary. 
Display of error bars is required. If the measurement errors are small, state so in the caption.

To understand better the basics of scientific graphing refer to \cite{Tufte01}.


\subsection{Requirements for tables}
As with graphs, an explanation of the table is required in a numbered caption. The caption is placed below the table and describes its content.

Specification of the physical units are mandatory.

\subsection{Requirements for references}
Common scientific sources are textbooks or scientific journal articles. References to the Wikipedia are not accepted as valid scientific reference sources. The sources are preferably written in the format found in PRL, (Physical Review Letters).  At a minimum they must be numbered and include author information, a citation to the reference source, and the relevant pages must be specified.

A convenient way to manage bibliography and references is provided by \TeX \  \texttt{bibtex} package.

\section{Discussion}

This section includes the following topics:
\begin{enumerate}
    \item Comparison with theory, including quantitative comparison using specific criterion.
    \item  Somewhere in the table (the column heading or the caption) should indicate the accuracy of the numerical values based on the relative error (e.g. $\frac{\delta x}{x}$).
    \item Discussion of errors, their types, sources and ways of reducing them; In particular, identifying the most significant contributor to the error. The discussion is based on the analysis of the significant errors (i.e., the quantitative basis for the arguments raised in the discussion).
\end{enumerate}


\section{Conclusions}
The conclusions are based on numerical values (as opposed to general statements). 

The conclusions are related to the questions raised in the \textbf{Summary} section.

The conclusions are relevant to the specific experiment carried out by the students.

The section does not exceed three paragraphs, which are written in narrative form as opposed to list.

\section*{Appendix}
Optional section. Contains formulae and discussion of significant errors and their implications. As required it may contain other derivations and explanations. 

Report might contain relevant snippets of code.
\begin{lstlisting}[language=Python, caption=Python code snippet example]
import numpy as np
import matplotlib.pyplot as plt

x = np.linspace(0, 2*np.pi)
y = np.sin(x)

plt.plot(x,y, '.')
plt.xlabel('x axis')
plt.ylabel('y axis')
plt.grid()
\end{lstlisting}


\printbibliography  % Print the bibliography
  

\end{document}